% SIAM Shared Information Template
% This is information that is shared between the main document and any
% supplement. If no supplement is required, then this information can
% be included directly in the main document.


% Packages and macros go here
\usepackage{lipsum}
\usepackage{amsfonts}
\usepackage{graphicx}
\usepackage{epstopdf}
\usepackage{algorithmic}

\usepackage[utf8]{inputenc}
\usepackage[T1]{fontenc}
\usepackage{nccmath}
% \usepackage{geometry}
\usepackage{latexsym}
\usepackage{amsmath}
\usepackage{graphics}
\usepackage{fullpage}
\usepackage{epsfig}
\usepackage{relsize}
\usepackage{amsfonts}
\usepackage{mathrsfs}
\usepackage{makeidx}
% \usepackage{comment}
\usepackage{latexsym}
\usepackage{dcolumn}
% \usepackage{tocloft} % spacing for tableofcontents
\usepackage{calc}
\usepackage{pgfplots}
\usepackage{pgfplotstable}
\usepackage{makecell}
\usepackage{mathtools}
\usepackage{xfrac}
\usepackage{url}
\usepackage{mathtools}
% \usepackage[
%      colorlinks = true,
%      citecolor = blue,
%      urlcolor = blue
%      ]{hyperref}

\usepackage{tikz}
\usetikzlibrary{arrows.meta, calc}
\usepackage[ocgcolorlinks]{ocgx2}
\usepackage[shortlabels]{enumitem}
% \excludecomment{codes}

\setlength{\abovecaptionskip}{0pt plus 0pt minus 2pt}
\ifpdf
  \DeclareGraphicsExtensions{.eps,.pdf,.png,.jpg}
\else
  \DeclareGraphicsExtensions{.eps}
\fi

% Add a serial/Oxford comma by default.
\newcommand{\creflastconjunction}{, and~}

% Used for creating new theorem and remark environments
\newsiamremark{remark}{Remark}
\newsiamremark{example}{Example}
\newsiamremark{hypothesis}{Hypothesis}
\crefname{hypothesis}{Hypothesis}{Hypotheses}
\newsiamthm{claim}{Claim}
% \newsiamthm{definition}{Definition}

% -------------------------------------------------
% {{{ Environments and some setups
% -------------------------------------------------
% \theoremstyle{plain}
% \newtheorem{theorem}{Theorem}[section]
% \newtheorem{corollary}[theorem]{Corollary}
% \newtheorem{lemma}[theorem]{Lemma}
% \newtheorem{proposition}[theorem]{Proposition}
%
% \theoremstyle{definition}
% \newtheorem{definition}[theorem]{Definition}
% \newtheorem{hypothesis}[theorem]{Hypothesis}
% \newtheorem{assumptions}[theorem]{Assumption}
%
% \newtheorem{remark}[theorem]{Remark}
% \newtheorem{conjecture}[theorem]{Conjecture}
% \newtheorem{example}[theorem]{Example}
% \newtheorem{notation}[theorem]{Notation}
%
% \numberwithin{equation}{section}
% \allowdisplaybreaks
% }}}


% -------------------------------------------------
% {{{ Biber related
% -------------------------------------------------
\usepackage[strict = true,
            style = english]{csquotes}
% \usepackage{mathscinet}
\DeclareTextCommand{\polhk}{OT4}{\k}
\DeclareTextCommand{\polhk}{T1}{\k}
\usepackage[backend = biber,
						% style = alphabetic,
						style = numeric,
						natbib = true,
            sorting = nyt,
            maxbibnames=99,
            abbreviate = true,
           ]{biblatex}
\AtEveryBibitem{% Clean up the bibtex rather than editing it
 \clearfield{doi}
 \clearfield{url}
 % \clearfield{issn}
 % \clearlist{location}
 \clearfield{month}
 \clearfield{series}

 \ifentrytype{book}{}{% Remove publisher and editor except for books
  \clearlist{publisher}
  \clearname{editor}
 }
}
\renewcommand*{\bibfont}{\fontsize{11}{11}}

% \addbibresource{All.bib}
\addbibresource{./draft_rand_int_2_biber.bib}
% }}}

% Sets running headers as well as PDF title and authors
% \headers{Invariant measure for SHE}{L. Chen and N. Eisenberg}

% Title. If the supplement option is on, then "Supplementary Material"
% is automatically inserted before the title.
\title{Invariant measures for the nonlinear stochastic heat equation with no drift term\thanks{Submitted to the editors Dec. 12, 2022.
\funding{This work (L.C.) was partially supported by the \textit{Travel Support for Mathematicians} from the \textit{Simons Foundation}.}}}

% Authors: full names plus addresses.
\author{
  Le Chen\thanks{Department of Mathematics and Statistics, Auburn University, Auburn, Alabama, USA.
  (L.C.: \email{le.chen@auburn.edu}, \url{http://webhome.auburn.edu/\~lzc0090};
   N.E.: \email{nze0019@auburn.edu} and \email{nickeisenberg@gmail.com}).}
  \and
  Nicholas Eisenberg\footnotemark[2]
}

\usepackage{amsopn}
\DeclareMathOperator{\diag}{diag}


%%% Local Variables:
%%% mode:latex
%%% TeX-master: "ex_article"
%%% End:
